\section{Introduction}
\label{sec:intro}

\subsection{The question}

Modern genome assembly reconstructs an unknown genome from many short
overlapping reads. Two lines of theory ask two different questions about this
problem. \emph{Information-feasibility} asks whether the reads contain enough
information to determine the genome; the answer depends on the repeat structure
of the genome and on the read length, and Shomorony, Kim, Courtade, and Tse
\cite{shomorony2016} gave repeat-bridging conditions under which a specific
algorithm recovers the true circular sequence up to cyclic shift. \emph{Maximum
likelihood} instead asks which genome is the most probable source of the
observed reads; Medvedev and Brudno \cite{medvedev2009} formulated this as an
optimization problem over candidate genomes and solved a relaxed version by
min-cost flow.

Shomorony et al.\ close their Discussion with a question that sits between the
two lines \cite[Discussion]{shomorony2016}:

\begin{quote}
``Understanding whether bridging conditions can be used to guarantee that the
maximum-likelihood sequence is the true sequence is currently an open
question.''
\end{quote}

This manuscript studies that question as a mathematical problem --- but not by
pretending it is a single well-posed theorem. The sentence names a likelihood
formulation only by citation and leaves several choices open. Our first
contribution is to make those choices explicit; our second is to report what is
provable or refutable for each resulting formulation.

\subsection{Why maximum likelihood is not automatic}

A naive argument says the question should be trivial. The reads were generated
from the true genome, so the true genome should maximize their likelihood. The
argument fails for a structural reason: the read distribution can be
\emph{non-identifiable}. The reads are drawn from the multiset of length-$L$
windows of the genome, so two genomes with the same window spectrum induce
exactly the same read distribution and always tie.

\begin{example}[Structural ambiguity before the repeat boundary]
\label{ex:ABAB}
Let the alphabet be $\{A,B\}$ and let $L=2$. The circular genomes
\[
  S = AAABAB, \qquad D = AABAAB
\]
of length $6$ are not cyclic rotations of one another, yet their circular
$2$-mer spectra agree:
\[
  \spec_2(S) = \spec_2(D) = \{\,AA:2,\ AB:2,\ BA:2\,\}.
\]
Indeed $S$ has windows $AA,AA,AB,BA,AB,BA$ and $D$ has windows
$AA,AB,BA,AA,AB,BA$. Under uniform-start error-free sampling the observed read
distribution is the same for both genomes, so \emph{no} likelihood objective
--- exact, approximate, or flow-based --- can prefer $S$ to $D$ here.
\statusline{\srcfact{} for the spectra and the rotation claim;
\mathfact{} for the resulting tie. This is the ``$AAABAB/AABAAB$''
example recorded in the project synthesis.}
\end{example}

The lesson of Example~\ref{ex:ABAB} is that maximum likelihood cannot manufacture
identifiability. Bridging conditions exist precisely to remove the structural
ambiguity: they are hypotheses on the repeat structure and read coverage of the
\emph{true} genome that, in Shomorony et al.'s theorem, force the true genome to
be the unique Eulerian object consistent with the reads. The published question
asks whether those same hypotheses also make the true genome optimal for a
likelihood objective. It is therefore a question about whether a
\emph{reconstruction} guarantee transfers to an \emph{optimization} guarantee.

\subsection{What the question does not say}

Before any theorem can be attributed to the published sentence, four modeling
axes must be fixed, and we will argue in Section~\ref{sec:open} that the
sentence fixes none of them conclusively. The named formulation of
\cite{medvedev2009} contains at least four distinct objects:

\begin{enumerate}[label=(\roman*)]
  \item the exact global read-count multinomial with candidate-intrinsic length;
  \item its separable product-of-binomial-marginals approximation, which replaces
        the candidate length by an external genome size;
  \item the Section~6.2 bidirected read-overlap-graph flow, whose feasible
        objects are flows and need not be sequences; and
  \item the underlying broad ``maximum-likelihood principle.''
\end{enumerate}

Independently, the sentence does not say whether reads are oriented strings or
reverse-complement molecule classes; whether candidate genomes are restricted to
the true length; or whether the conclusion asserts that the truth is \emph{some}
maximizer or that \emph{every} maximizer is the truth up to a genome equivalence.
These axes matter, and we treat them as part of the mathematical object rather
than as notation.

\subsection{Contributions}

This manuscript does four things.

\paragraph{Negative finite results.}
We exhibit kernel-checked finite counterexamples showing that the
information-feasibility hypothesis does \emph{not} force the truth to be a
likelihood maximizer under several precise readings of the 2016 question: the
exact multinomial over all circular candidates (Theorem~\ref{thm:cex-exact}), the
fixed-length binomial approximation over all circular candidates
(Theorem~\ref{thm:cex-binomial}), and the reverse-complement molecule reading of
the Section~6.2 flow (Theorem~\ref{thm:cex-molecule}). A candidate-set inclusion
lemma (Lemma~\ref{lem:negative-transfer}) shows that a same-length refutation
transfers to every superclass of candidates, so these results are not confined
to length-restricted comparisons.

\paragraph{A positive same-length rigidity theorem.}
Under \emph{strict oriented} read types, Shomorony's $\Is$, and the
same-length spelled-candidate restriction, the truth's spectrum is the unique
positive circulation of its window-support graph (Theorem~\ref{thm:rigidity}).
Consequently every same-length Section~6.2 spelled candidate ties the truth under
any spectrum-based objective, and no strict same-length counterexample exists for
\emph{any} genome length, read length, or alphabet. The equal-length restriction
is essential: allowing one extra symbol reintroduces a strict witness
(Example~\ref{ex:variable-length}).

\paragraph{A repaired population model.}
We present a new model, motivated by the finite failures rather than by the
historical papers: replace empirical read counts by the exact population read
distribution. A Gibbs/KL identity (Lemma~\ref{lem:gibbs}) makes the truth a
population maximizer \emph{unconditionally}. Uniqueness then reduces to spectrum
identifiability, and for the project's oriented, primitive, P2-admissible
candidate class this is supplied by a complete-spectrum uniqueness theorem of
Bresler, Bresler, and Tse \cite{bresler2013}. The resulting
Theorem~\ref{thm:population} is the current positive high point of the project.
It does not settle the finite 2016 question, and we state the gap explicitly.

\paragraph{Exposition as a first-class concern.}
Because the mathematical content is an interface between under-specified
sources and project-level repairs, correctness includes keeping the epistemic
layers separate. Theorem statements name their assumption surface; every
nontrivial claim carries a status tag (\statuslegend); and
Appendices~\ref{app:modeling} and~\ref{app:ledger} explain the modeling choices
and consolidate the status of each claim.

\subsection{How to read this paper}

Section~\ref{sec:model} fixes notation and the source-faithful sequencing and
bridging model. Section~\ref{sec:open} states the published question and
determines what the sources do and do not fix. Section~\ref{sec:finite} reports
the results about the literature-derived finite formulations.
Section~\ref{sec:population} introduces the population model and its uniqueness
theorem. Section~\ref{sec:discussion} draws the two failure modes together and
lists what remains open. Appendix~\ref{app:modeling} explains the modeling
choices deliberately; Appendix~\ref{app:ledger} is the consolidated epistemic
ledger; Appendix~\ref{app:reproduction} records how to rebuild the paper and
re-run every computational certificate.

A reader who wants only the mathematical story can read
Sections~\ref{sec:model}--\ref{sec:discussion} and treat the appendices as
reference. A reader who wants to scrutinize the source correspondence should
read Appendix~\ref{app:modeling} first.
